% ===========================================================================
%  IEEE Conference Paper
%  Zero-Sum Robust Deep RL with Transformer-Based Disturbance Detection
%  and Policy Mixing for Real-Time Control on Embedded GPUs
% ===========================================================================
\documentclass[conference]{IEEEtran}

% ---- packages -------------------------------------------------------------
\usepackage{amsmath}
\usepackage{amssymb}
\usepackage{algorithmic}
\usepackage{algorithm}
\usepackage{graphicx}
\usepackage{booktabs}
\usepackage{hyperref}
\usepackage{xcolor}
\usepackage{multirow}

% ---- hyperref setup --------------------------------------------------------
\hypersetup{
  colorlinks=true,
  linkcolor=black,
  citecolor=blue!60!black,
  urlcolor=blue!60!black
}

% ---- math helpers ----------------------------------------------------------
\newcommand{\E}{\mathbb{E}}
\newcommand{\R}{\mathbb{R}}
\newcommand{\cS}{\mathcal{S}}
\newcommand{\cA}{\mathcal{A}}
\newcommand{\cP}{\mathcal{P}}
\newcommand{\cR}{\mathcal{R}}
\newcommand{\cL}{\mathcal{L}}
\newcommand{\cD}{\mathcal{D}}
\newcommand{\cT}{\mathcal{T}}
\newcommand{\piopt}{\pi^{\text{opt}}}
\newcommand{\pirob}{\pi^{\text{rob}}}
\newcommand{\piadv}{\pi^{\text{adv}}}

% ===========================================================================
\title{Zero-Sum Robust Deep Reinforcement Learning with\\
Transformer-Based Disturbance Detection and Policy Mixing\\
for Real-Time Control on Embedded GPUs}

\author{
\IEEEauthorblockN{Author A\IEEEauthorrefmark{1},
Author B\IEEEauthorrefmark{1},
Author C\IEEEauthorrefmark{2},
Author D\IEEEauthorrefmark{1}}
\IEEEauthorblockA{\IEEEauthorrefmark{1}Department of Computer Science,
University of Placeholder, City, Country\\
Email: \{author.a, author.b, author.d\}@placeholder.edu}
\IEEEauthorblockA{\IEEEauthorrefmark{2}Department of Electrical Engineering,
Institute of Technology, City, Country\\
Email: author.c@placeholder.edu}
}

\begin{document}
\maketitle

\begin{IEEEkeywords}
robust reinforcement learning, zero-sum Markov games, adversarial training,
transformer anomaly detection, policy mixing, TensorRT, embedded deployment
\end{IEEEkeywords}

% ===========================================================================
\begin{abstract}
% ===========================================================================
Deploying reinforcement learning (RL) policies on physical robots requires
robustness to unmodeled disturbances, parameter uncertainty, and the
sim-to-real gap.  We present a methodology for robust control that frames
the problem as a two-player zero-sum Markov game in which a learned
adversary systematically probes the weaknesses of the control policy during
training.  Our architecture maintains two distinct policies: an
\emph{optimal} policy $\piopt$ trained on the nominal environment to
maximize expected return, and a \emph{robust} policy $\pirob$ co-trained
against the adversary to maximize worst-case return.  At runtime, a
lightweight transformer encoder processes a sliding window of recent
observations to produce a real-time disturbance estimate---both a detection
probability and an estimated magnitude.  A learned gating function converts
this estimate into a continuous mixing coefficient $\alpha_t \in [0,1]$
that blends the two policies, yielding actions
$a_t = \alpha_t\,\pirob(s_t) + (1-\alpha_t)\,\piopt(s_t)$.  The training
pipeline is implemented in PyTorch with Gymnasium and MuJoCo, with a
planned migration to JAX and MJX for GPU-vectorized environment
parallelism.  For deployment, all neural network components are exported
through ONNX to TensorRT engines optimized for the NVIDIA Jetson platform,
and integrated into a ROS~2 control node that achieves sub-two-millisecond
inference latency at up to 500\,Hz.  We provide a complete data-logging
specification, three detailed algorithmic procedures covering adversarial
training, transformer detector training, and real-time inference, and a
comprehensive evaluation plan with ablation studies and baselines.  This
paper describes the full methodology and system design; empirical results
are deferred to a follow-up study.
\end{abstract}

% ===========================================================================
\section{Introduction}\label{sec:intro}
% ===========================================================================

Reinforcement learning has achieved impressive results in simulated
continuous-control benchmarks, yet transferring learned policies to physical
robots remains a significant challenge.  Real-world systems are subject to
disturbances---external forces, friction changes, payload variation,
sensor noise, and actuator degradation---that are difficult to capture
faithfully in a simulator.  The resulting \emph{sim-to-real gap} frequently
causes policies that perform well in simulation to fail when deployed on
hardware~\cite{zhao2020simtoreal}.

A principled approach to this problem is \emph{robust} reinforcement
learning, which seeks policies that perform well not only under nominal
conditions but also under worst-case perturbations.  One particularly
elegant formulation models the interaction between the controller and the
environment disturbance as a \emph{two-player zero-sum Markov
game}~\cite{littman1994markov}: the controller maximizes cumulative reward
while a learned adversary minimizes it, and training converges toward a
Nash equilibrium that provides formal worst-case guarantees.  This approach,
pioneered by Robust Adversarial Reinforcement Learning
(RARL)~\cite{pinto2017robust} and the $H_\infty$ control
perspective~\cite{morimoto2005robust}, has since been extended to
state-adversarial formulations~\cite{zhang2020robust} and adversarial
populations~\cite{vinitsky2020robust}.

Despite these advances, several gaps remain in the literature.  First,
existing adversarial training methods produce a \emph{single} policy that
compromises between nominal performance and worst-case robustness; no
mechanism exists to recover full nominal performance when disturbances are
absent.  Second, at deployment time, the policy has no explicit
\emph{awareness} of whether or not a disturbance is currently active: it
applies the same conservative strategy regardless of conditions.  Third,
transformer-based anomaly detection has shown strong results in multivariate
time-series domains~\cite{tuli2022tranad, xu2022anomaly}, yet this
capability has not been leveraged for runtime disturbance detection in RL
control systems.  Finally, few works address the complete pipeline from
adversarial training through model export to real-time inference on embedded
GPU platforms with ROS~2 integration.

This paper addresses these gaps with the following contributions:
\begin{itemize}
\item A \textbf{dual-policy mixing architecture} that maintains separate
      optimal and robust policies and blends their actions through a
      continuous gating coefficient, preserving nominal performance while
      providing worst-case robustness on demand.
\item A \textbf{transformer-based disturbance detector} that operates on a
      sliding window of observations and outputs a detection probability
      and estimated disturbance magnitude in real time, providing the signal
      needed for the gating decision.
\item A \textbf{unified training pipeline} with a detailed data-logging
      specification that supports both off-policy and on-policy adversarial
      training algorithms and generates the supervised dataset for the
      transformer detector.
\item A \textbf{deployment pipeline} that exports all neural network
      components to ONNX, optimizes them with TensorRT for the NVIDIA
      Jetson platform, and integrates them into a ROS~2 node meeting
      sub-two-millisecond latency budgets.
\item A \textbf{comprehensive evaluation plan} including ablation studies,
      baselines, and safety checks, designed to rigorously validate the
      approach once empirical experiments are conducted.
\end{itemize}

The remainder of this paper is organized as follows.
Section~\ref{sec:background} reviews the necessary technical background.
Section~\ref{sec:related} surveys related work.
Section~\ref{sec:formulation} formalizes the problem.
Section~\ref{sec:method} presents the method in detail, including
algorithms and data specifications.
Section~\ref{sec:implementation} describes the system implementation.
Section~\ref{sec:evaluation} outlines the evaluation plan.
Section~\ref{sec:discussion} discusses limitations, and
Section~\ref{sec:conclusion} concludes.

% ===========================================================================
\section{Background and Preliminaries}\label{sec:background}
% ===========================================================================

% ---------------------------------------------------------------------------
\subsection{Markov Decision Processes}\label{sec:bg_mdp}
% ---------------------------------------------------------------------------

A Markov Decision Process (MDP) is defined by the tuple
$(\cS, \cA, P, R, \gamma)$, where $\cS$ is the state space,
$\cA$ is the action space, $P: \cS \times \cA \times \cS \to [0,1]$ is
the transition kernel, $R: \cS \times \cA \to \R$ is the reward function,
and $\gamma \in [0,1)$ is the discount factor.  A stochastic policy
$\pi: \cS \to \Delta(\cA)$ maps states to distributions over actions.  The
state-value function and action-value function under policy $\pi$ are
\begin{align}
V^\pi(s) &= \E_\pi\!\left[\sum_{t=0}^{\infty}\gamma^t R(s_t, a_t)
  \;\middle|\; s_0 = s\right], \label{eq:value}\\
Q^\pi(s,a) &= \E_\pi\!\left[\sum_{t=0}^{\infty}\gamma^t R(s_t, a_t)
  \;\middle|\; s_0 = s, a_0 = a\right]. \label{eq:qvalue}
\end{align}
These satisfy the Bellman equation
\begin{equation}\label{eq:bellman}
V^\pi(s) = \E_{a\sim\pi(\cdot|s)}\!\Big[R(s,a) +
  \gamma\,\E_{s'\sim P(\cdot|s,a)}V^\pi(s')\Big].
\end{equation}
The goal of reinforcement learning is to find a policy $\pi^*$ that
maximizes $V^\pi(s)$ for all $s \in \cS$.

% ---------------------------------------------------------------------------
\subsection{Two-Player Zero-Sum Markov Games}\label{sec:bg_game}
% ---------------------------------------------------------------------------

A two-player zero-sum Markov game~\cite{littman1994markov} extends the MDP
to $(\cS, \cA_1, \cA_2, P, R, \gamma)$, where $\cA_1$ and $\cA_2$ are
the action spaces of Player~1 (the controller) and Player~2 (the
adversary/disturbance), respectively.  The transition kernel becomes
$P: \cS \times \cA_1 \times \cA_2 \times \cS \to [0,1]$, and the reward
$R(s, a_1, a_2)$ is received by the controller while the adversary
receives $-R(s, a_1, a_2)$.  A pair of policies
$(\pi_1^*, \pi_2^*)$ constitutes a \emph{Nash equilibrium} if
\begin{equation}\label{eq:nash}
V^{\pi_1^*,\pi_2}(s) \;\geq\; V^{\pi_1^*,\pi_2^*}(s) \;\geq\;
V^{\pi_1,\pi_2^*}(s)
\end{equation}
for all $s$, $\pi_1$, $\pi_2$.  The \emph{minimax value} is
\begin{equation}\label{eq:minimax}
V^*(s) = \max_{\pi_1}\min_{\pi_2}\;
  \E\!\left[\sum_{t=0}^{\infty}\gamma^t R(s_t, a_t^1, a_t^2)
  \;\middle|\; s_0=s\right].
\end{equation}
This formulation provides a natural framework for robust control: the
controller learns to maximize reward against the strongest possible
adversary, yielding a policy with formal worst-case performance
guarantees.

% ---------------------------------------------------------------------------
\subsection{Deep Reinforcement Learning Algorithms}\label{sec:bg_drl}
% ---------------------------------------------------------------------------

We consider four deep RL algorithms as candidate training methods for
both the controller and the adversary.

\textbf{DDPG}~\cite{lillicrap2016continuous} is an off-policy actor-critic
algorithm for continuous action spaces.  It maintains a deterministic
policy $\mu_\theta(s)$, a critic $Q_\phi(s,a)$, a replay buffer for
experience storage, and target networks $\mu_{\theta'}$, $Q_{\phi'}$
updated via Polyak averaging.  The critic is trained by minimizing the
temporal-difference (TD) error, and the actor is updated by applying the
deterministic policy gradient~\cite{silver2014deterministic}.

\textbf{TD3}~\cite{fujimoto2018addressing} improves upon DDPG with three
modifications: (i)~twin critics $Q_{\phi_1}$, $Q_{\phi_2}$ with the
minimum taken as the target (reducing overestimation bias), (ii)~delayed
policy updates (the actor is updated less frequently than the critic), and
(iii)~target policy smoothing (noise is added to the target action to
regularize the critic).

\textbf{SAC}~\cite{haarnoja2018soft} is an off-policy algorithm based on
the maximum-entropy RL framework.  The objective augments the expected
return with an entropy term:
\begin{equation}\label{eq:sac_obj}
J_{\text{SAC}}(\pi) = \E_\pi\!\left[\sum_{t=0}^{\infty}\gamma^t
  \Big(R(s_t,a_t) + \beta\,\mathcal{H}\!\big(\pi(\cdot|s_t)\big)\Big)
  \right],
\end{equation}
where $\beta > 0$ is the temperature parameter (often automatically
tuned) and $\mathcal{H}$ denotes the Shannon entropy.  SAC learns a
stochastic policy parameterized via the reparameterization trick, combined
with twin critics as in TD3.

\textbf{PPO}~\cite{schulman2017proximal} is an on-policy algorithm that
optimizes a clipped surrogate objective:
\begin{equation}\label{eq:ppo_obj}
J_{\text{PPO}}(\theta) = \E_t\!\left[\min\!\big(r_t(\theta)\hat{A}_t,\;
  \text{clip}(r_t(\theta), 1{-}\epsilon, 1{+}\epsilon)\hat{A}_t\big)
  \right],
\end{equation}
where $r_t(\theta) = \pi_\theta(a_t|s_t)/\pi_{\theta_\text{old}}(a_t|s_t)$
is the probability ratio and $\hat{A}_t$ is the generalized advantage
estimate (GAE).  PPO collects rollouts in batches and performs multiple
epochs of mini-batch optimization before discarding the data.

% ---------------------------------------------------------------------------
\subsection{Robust Reinforcement Learning Formulations}\label{sec:bg_robust}
% ---------------------------------------------------------------------------

Several formulations address robustness in RL.
\emph{Robust MDPs}~\cite{iyengar2005robust, nilim2005robust} replace the
single transition kernel $P$ with an \emph{uncertainty set}
$\cP \ni P$, and the agent optimizes against the worst-case element:
$V^*(s) = \max_\pi \min_{P \in \cP} V^\pi_P(s)$.  The
\emph{$H_\infty$ approach}~\cite{morimoto2005robust} penalizes the
sensitivity of the value function to disturbances, yielding a minimax
objective similar to~\eqref{eq:minimax} but with a continuous disturbance
signal rather than a discrete adversary.

\emph{Adversarial training} (RARL)~\cite{pinto2017robust} operationalizes
the zero-sum game by training a neural-network adversary that applies
forces to the agent's body, alternating between controller and adversary
updates.  This has been extended to \emph{state-adversarial}
MDPs~\cite{zhang2020robust}, where the adversary perturbs the
observations rather than the dynamics, and \emph{action-robust}
formulations~\cite{tessler2019action}, where the adversary modifies the
executed action.  A comprehensive survey of robust RL methods is provided
by Moos~et~al.~\cite{moos2022robust}.

% ===========================================================================
\section{Related Work}\label{sec:related}
% ===========================================================================

% ---------------------------------------------------------------------------
\subsection{Adversarial and Robust Reinforcement Learning}
% ---------------------------------------------------------------------------

The idea of training RL agents against adversarial perturbations
originates with Morimoto and Doya~\cite{morimoto2005robust}, who cast
robust control as a differential game and solved it with actor-critic
methods.  Pinto~et~al.~\cite{pinto2017robust} introduced RARL, in which
a destabilizing adversary applies forces to the agent's body during
training, demonstrating improved transfer to environments with
modified physical parameters.  Pan~et~al.~\cite{pan2019risk} proposed a
risk-averse variant of RARL that incorporates CVaR-based objectives to
further reduce tail risk.  Vinitsky~et~al.~\cite{vinitsky2020robust}
showed that training against an \emph{ensemble} of adversarial policies,
rather than a single adversary, prevents the controller from overfitting
to one attack strategy and improves generalization.

Zhang~et~al.~\cite{zhang2020robust} formulated the state-adversarial MDP
(SA-MDP), in which the adversary perturbs the agent's observations
within a bounded $\ell_p$ ball.  They derived a regularizer that
encourages policy smoothness with respect to state perturbations.
Oikarinen~et~al.~\cite{oikarinen2021robust} extended this idea by
incorporating a certified adversarial loss during training, providing
provable robustness bounds for linear function approximators.
Sun~et~al.~\cite{sun2022strongest} studied the problem of finding
\emph{optimal} adversarial attacks against deep RL agents, characterizing
the strongest possible attack under budget constraints and showing that
many existing attack heuristics are suboptimal.

A comprehensive survey by Moos~et~al.~\cite{moos2022robust} categorizes
robust RL methods along several axes: uncertainty sets, adversary types,
and solution concepts.  A common observation across this body of work is
that adversarial training is performed entirely at \emph{training time}:
the resulting policy applies the same robust strategy at every timestep
during deployment, regardless of whether a disturbance is actually
present.  This conflates nominal and adversarial operation, sacrificing
performance when conditions are benign.  Our work addresses this gap by
introducing a runtime disturbance detection mechanism that gates between
nominal-optimal and adversarially-robust policies.

% ---------------------------------------------------------------------------
\subsection{Transformers in Time-Series and Reinforcement Learning}
% ---------------------------------------------------------------------------

The transformer architecture~\cite{vaswani2017attention}, originally
developed for natural language processing, has been successfully applied
to time-series analysis and anomaly detection.
TranAD~\cite{tuli2022tranad} uses an adversarially-trained
transformer with attention-based reconstruction for detecting anomalies
in multivariate sensor streams, achieving state-of-the-art results on
industrial monitoring benchmarks.  The Anomaly
Transformer~\cite{xu2022anomaly} introduces an \emph{association
discrepancy} mechanism that distinguishes anomalous timesteps from normal
ones by comparing learned associations with a prior distribution.

In the RL domain, Chen~et~al.~\cite{chen2021decision} proposed the
Decision Transformer, which formulates RL as a sequence-modeling problem
and conditions action generation on desired return-to-go, states, and
actions.  While this work demonstrates the power of transformers for
representing sequential decision-making, it does not address disturbance
detection.

To our knowledge, no prior work uses transformers specifically for
\emph{runtime disturbance detection} in the context of robust RL,
nor does any work use such a detector to gate between separate
optimal and robust control policies.

% ---------------------------------------------------------------------------
\subsection{Sim-to-Real Transfer and Deployment}
% ---------------------------------------------------------------------------

Sim-to-real transfer remains a central challenge in robot learning.
\emph{Domain randomization}~\cite{tobin2017domain} trains policies on a
distribution of environment parameters (visual appearance, physical
properties) so that the real world appears as just another sample from the
training distribution.  Peng~et~al.~\cite{peng2018simtoreal}
demonstrated that randomizing dynamics parameters during training---mass,
friction, damping---enables zero-shot transfer of manipulation policies.
Zhao~et~al.~\cite{zhao2020simtoreal} provide a survey of sim-to-real
transfer methods, covering domain randomization, system identification,
and progressive adaptation.

Lee~et~al.~\cite{lee2020learning} used a teacher-student framework for
quadrupedal locomotion: a privileged teacher trained in simulation
distills its knowledge into a student that relies only on proprioceptive
observations available on the real robot.  This approach has become a
standard paradigm in legged locomotion.

On the simulation side, MuJoCo~\cite{todorov2012mujoco} remains the
dominant physics engine for contact-rich control.  The recent
MJX extension~\cite{zakka2025mujoco} provides a JAX-based implementation
of MuJoCo that enables massively parallel environment simulation on GPUs.
Brax~\cite{freeman2021brax} offers a similar JAX-based differentiable
physics engine.  Gymnasium~\cite{towers2024gymnasium} provides the
standard Python interface for RL environments.

% ---------------------------------------------------------------------------
\subsection{Embedded Inference for Real-Time Control}
% ---------------------------------------------------------------------------

Deploying neural-network policies on embedded platforms requires careful
attention to inference latency and throughput.  NVIDIA's
TensorRT~\cite{nvidia2024tensorrt} is a high-performance inference
optimizer that applies layer fusion, kernel auto-tuning, precision
calibration (FP16, INT8), and memory optimization to minimize latency on
NVIDIA GPUs.  Jeon~et~al.~\cite{jeon2022tensorrt} developed a framework
and optimization methodology specifically for TensorRT inference on Jetson
boards, achieving significant speedups over native PyTorch execution.

The ONNX ecosystem~\cite{onnxruntime2022transformer} provides a standard
interchange format for neural networks, enabling models trained in any
framework (PyTorch, TensorFlow, JAX) to be converted to a common
representation and subsequently optimized by TensorRT.  This is
particularly important for transformer architectures, where INT8
quantization of attention layers requires careful calibration.

NVIDIA's Isaac ROS~\cite{nvidia2024isaacros} provides GPU-accelerated
ROS~2 packages built on the NITROS framework, which enables zero-copy
GPU tensor transfer between ROS nodes.  This infrastructure is critical
for achieving the low latencies required by real-time control.

Despite the maturity of individual components, few works address the
\emph{complete pipeline} from adversarial RL training in simulation through
ONNX export and TensorRT optimization to real-time deployment on Jetson
hardware with ROS~2 integration.  Our work provides a detailed design for
this end-to-end pipeline.

% ---------------------------------------------------------------------------
\subsection{Comparison to Our Approach}
% ---------------------------------------------------------------------------

Table~\ref{tab:comparison} summarizes how our approach relates to prior
work.  The key differentiators are: (i)~the dual-policy architecture
with continuous gating, (ii)~the use of a transformer for runtime
disturbance detection to drive the gating signal, and (iii)~the
end-to-end deployment pipeline targeting embedded GPUs with real-time
ROS~2 integration.  No prior work combines all three elements.

\begin{table}[t]
\centering
\caption{Feature comparison with related approaches.}
\label{tab:comparison}
\renewcommand{\arraystretch}{1.15}
\small
\begin{tabular}{@{}lcccc@{}}
\toprule
\textbf{Feature} & \textbf{RARL} & \textbf{SA-MDP} &
  \textbf{DR} & \textbf{Ours} \\
\midrule
Adversarial training      & \checkmark & \checkmark & --         & \checkmark \\
Dual-policy mixing        & --         & --         & --         & \checkmark \\
Runtime detection         & --         & --         & --         & \checkmark \\
Transformer detector      & --         & --         & --         & \checkmark \\
TensorRT deployment       & --         & --         & --         & \checkmark \\
ROS~2 integration         & --         & --         & --         & \checkmark \\
\bottomrule
\end{tabular}
\vspace{-0.3cm}
\end{table}

% ===========================================================================
\section{Problem Formulation}\label{sec:formulation}
% ===========================================================================

% ---------------------------------------------------------------------------
\subsection{Two-Player Zero-Sum Markov Game}
% ---------------------------------------------------------------------------

We consider a two-player zero-sum Markov game defined by the tuple
$(\cS, \cA^c, \cA^d, P, R, \gamma)$, where $\cS \subseteq \R^n$ is the
continuous state space, $\cA^c \subseteq \R^{m_c}$ is the controller's
action space, and $\cA^d \subseteq \R^{m_d}$ is the disturbance
(adversary) action space.  The transition dynamics are governed by
$P(s' \mid s, a^c, a^d)$, where $s \in \cS$, $a^c \in \cA^c$,
$a^d \in \cA^d$, and $s' \in \cS$.  The scalar reward
$r(s, a^c, a^d) \in \R$ is received by the controller; the adversary
receives $-r(s, a^c, a^d)$.

The controller seeks a policy $\pi^c: \cS \to \Delta(\cA^c)$ that
maximizes the expected discounted return, while the adversary seeks a
policy $\pi^d: \cS \to \Delta(\cA^d)$ that minimizes it.  A Nash
equilibrium $(\pi^{c,*}, \pi^{d,*})$ satisfies
\begin{equation}\label{eq:nash_formal}
V^{\pi^{c,*}, \pi^d}(s) \geq V^{\pi^{c,*}, \pi^{d,*}}(s) \geq
V^{\pi^c, \pi^{d,*}}(s)
\end{equation}
for all $s \in \cS$, $\pi^c$, $\pi^d$.  The minimax value function is
\begin{equation}\label{eq:minimax_formal}
V^*(s) = \max_{\pi^c} \min_{\pi^d} \;\E\!\left[
  \sum_{t=0}^{\infty} \gamma^t\, r(s_t, a_t^c, a_t^d)
  \;\middle|\; s_0 = s\right].
\end{equation}

% ---------------------------------------------------------------------------
\subsection{Disturbance Model}\label{sec:disturbance_model}
% ---------------------------------------------------------------------------

The adversary's action $a^d \in \cA^d$ represents an external disturbance
applied to the system.  We consider two concrete instantiations:
\begin{enumerate}
\item \textbf{Additive force/torque perturbation.}  The adversary applies
  a force or torque vector to one or more joints or links of the robot.
  The perturbed dynamics become
  $\dot{q} = f(q, \dot{q}, a^c) + B\,a^d$, where
  $B \in \R^{n_q \times m_d}$ is a known input matrix selecting the
  affected degrees of freedom.
\item \textbf{External body force.}  The adversary applies an external
  wrench (force and torque) to a specified body, modeling pushes, wind,
  or contact with unmodeled objects.
\end{enumerate}

In both cases, the adversary's action is bounded:
\begin{equation}\label{eq:dist_budget}
\|a^d\|_2 \leq \epsilon,
\end{equation}
where $\epsilon > 0$ is the \emph{disturbance budget}.  The adversary
learns to choose the direction and magnitude of the perturbation within
this budget to maximally degrade the controller's performance.  The
disturbance can be further parameterized by a magnitude
$\delta \in [0, \delta_{\max}]$ and a unit direction vector
$\boldsymbol{d} \in \R^{m_d}$, $\|\boldsymbol{d}\|_2 = 1$, so that
$a^d = \delta\,\boldsymbol{d}$.

% ---------------------------------------------------------------------------
\subsection{Objective: Dual-Policy Robust Control}\label{sec:objective}
% ---------------------------------------------------------------------------

Rather than training a single policy that compromises between nominal
performance and worst-case robustness, we maintain two separate policies:
\begin{itemize}
\item $\piopt$: the \emph{optimal} policy, trained on the nominal
  environment (no adversary) to maximize $\E[\sum_t \gamma^t r_t]$.
\item $\pirob$: the \emph{robust} policy, co-trained with the adversary
  $\piadv$ in the zero-sum game to maximize
  $\min_{\pi^d} \E[\sum_t \gamma^t r_t]$.
\end{itemize}
Both policies share the same observation and action spaces.  At runtime,
the executed action is a convex combination:
\begin{equation}\label{eq:mixing}
a_t = \alpha_t \cdot \pirob(s_t) + (1 - \alpha_t) \cdot \piopt(s_t),
\end{equation}
where the mixing coefficient $\alpha_t \in [0,1]$ is produced by a
\emph{gating network} $g_\theta$ operating on a history of observations:
\begin{equation}\label{eq:gating}
\alpha_t = g_\theta(h_t), \quad
h_t = (o_{t-L+1},\; o_{t-L+2},\; \ldots,\; o_t),
\end{equation}
with $o_i \in \R^d$ denoting the observation at timestep $i$ and $L$
denoting the history window length.

When $\alpha_t \approx 0$, the system executes the nominal-optimal policy;
when $\alpha_t \approx 1$, it executes the robust policy.  This design
preserves the full nominal performance of $\piopt$ in undisturbed
conditions while providing the worst-case guarantees of $\pirob$ when
disturbances are detected.

% ---------------------------------------------------------------------------
\subsection{Real-Time Constraints}\label{sec:rt_constraints}
% ---------------------------------------------------------------------------

Real-time control imposes strict latency requirements.  Let
$\tau_{\text{budget}}$ denote the maximum allowable inference time per
control step.  For a control frequency of $f$ Hz, we require
\begin{equation}\label{eq:latency}
\tau_{\text{gate}} + \tau_{\text{policy}} + \tau_{\text{pub}} \leq
\tau_{\text{budget}} = \frac{1}{f}.
\end{equation}
Here $\tau_{\text{gate}}$ is the time for the transformer forward pass and
gating computation, $\tau_{\text{policy}}$ is the time for the two policy
forward passes (which may be pipelined), and $\tau_{\text{pub}}$ is the
time for ROS message serialization and publication.  For the target range
of $f \in [200, 500]$ Hz, this yields
$\tau_{\text{budget}} \in [2, 5]$ ms.  All neural network components must
be optimized to meet this budget on the NVIDIA Jetson platform.

% ===========================================================================
\section{Method}\label{sec:method}
% ===========================================================================

This section presents the proposed approach in detail.  We describe the
dual-policy architecture, the adversary training objective, the
transformer disturbance detector, the gating rule, the data-logging
specification, and the complete algorithmic procedures.

% ---------------------------------------------------------------------------
\subsection{Dual-Policy Architecture}\label{sec:dual_policy}
% ---------------------------------------------------------------------------

The optimal policy $\piopt$ and the robust policy $\pirob$ are trained
independently.  Both are parameterized as feedforward neural networks
(typically two hidden layers of 256 units with ReLU activations) that
map observations $o \in \R^d$ to actions $a \in \R^{m_c}$.

\textbf{Optimal policy training.}  The policy $\piopt$ is trained using a
standard deep RL algorithm (SAC, TD3, or PPO) on the nominal environment,
i.e., with the adversary disabled ($a^d = \boldsymbol{0}$).  This policy
maximizes the expected return under nominal dynamics and represents the
best achievable performance when no disturbance is present.

\textbf{Robust policy training.}  The policy $\pirob$ is trained in the
zero-sum game against a co-learned adversary $\piadv$.  The controller and
adversary are updated in alternation, as described in
Section~\ref{sec:adv_training}.  Upon convergence, $\pirob$ approximates
the maximin strategy and provides robust performance under the worst-case
disturbance within the budget $\epsilon$.

\textbf{Rationale for separate training.}  A single policy trained with
adversarial perturbations necessarily trades off nominal performance for
robustness: it learns conservative actions even in the absence of
disturbances.  By maintaining two separate policies, our architecture
avoids this trade-off.  The gating mechanism (Section~\ref{sec:gating})
selects between them based on the detected disturbance level, achieving
near-optimal nominal performance and near-optimal worst-case performance
simultaneously.

% ---------------------------------------------------------------------------
\subsection{Adversary Training Objective}\label{sec:adv_training}
% ---------------------------------------------------------------------------

The adversary $\piadv$ is a neural network with the same architecture as
the controller policies, mapping observations $o \in \R^d$ to disturbance
actions $a^d \in \R^{m_d}$.  The output layer uses a $\tanh$ activation
scaled by the disturbance budget $\epsilon$:
\begin{equation}\label{eq:adv_output}
a^d = \epsilon \cdot \tanh\!\big(\text{NN}_{\psi}(o)\big),
\end{equation}
which ensures $\|a^d\|_\infty \leq \epsilon$ by construction.

The adversary's objective is to minimize the controller's return, i.e., to
maximize the negated return:
\begin{equation}\label{eq:adv_obj}
J^{\text{adv}}(\piadv) = -\,J^c(\pirob, \piadv) =
-\,\E_{\pirob,\piadv}\!\left[\sum_{t=0}^{\infty}\gamma^t
  r(s_t, a_t^c, a_t^d)\right].
\end{equation}

Training proceeds by \emph{alternating optimization}:
\begin{enumerate}
\item Fix $\piadv$; update $\pirob$ to maximize $J^c(\pirob, \piadv)$
  using the chosen RL algorithm.
\item Fix $\pirob$; update $\piadv$ to maximize $J^{\text{adv}}(\piadv)$
  using the same (or a separate) RL algorithm.
\end{enumerate}
This alternation is related to \emph{fictitious play} in game theory
and to the \emph{double oracle} method for solving zero-sum games.  In
practice, both players share a single replay buffer (for off-policy
methods) and are updated with a fixed ratio of gradient steps per
environment step.

To prevent the adversary from becoming too strong too quickly
(which can destabilize the controller's learning), we employ curriculum
scheduling of the disturbance budget: $\epsilon$ is linearly increased
from $0$ to $\epsilon_{\max}$ over the first $N_{\text{warmup}}$ episodes.

% ---------------------------------------------------------------------------
\subsection{Transformer Disturbance Detector}\label{sec:transformer}
% ---------------------------------------------------------------------------

The disturbance detector is a lightweight transformer encoder that
processes a sliding window of recent observations and produces two
outputs: a disturbance detection probability $p_t \in [0,1]$ and an
estimated disturbance magnitude $\hat{\delta}_t \geq 0$.

\textbf{Input representation.}  The input to the transformer is a sequence
of $L$ observation vectors $h_t = (o_{t-L+1}, \ldots, o_t)$, where
each $o_i \in \R^d$.  The sequence is treated as a matrix
$\boldsymbol{H}_t \in \R^{L \times d}$.

\textbf{Positional encoding.}  Since the transformer architecture is
permutation-invariant, we add sinusoidal positional encodings to each
element of the sequence:
\begin{align}
\text{PE}(k, 2j)   &= \sin\!\left(\frac{k}{10000^{2j/d_{\text{model}}}}\right), \label{eq:pe_sin}\\
\text{PE}(k, 2j+1) &= \cos\!\left(\frac{k}{10000^{2j/d_{\text{model}}}}\right), \label{eq:pe_cos}
\end{align}
where $k \in \{0, \ldots, L-1\}$ is the position index and
$j \in \{0, \ldots, d_{\text{model}}/2 - 1\}$ indexes the dimension.

\textbf{Architecture.}  The encoder consists of $N_L \in \{2, 3, 4\}$
transformer layers, each containing multi-head self-attention with
$N_H \in \{2, 4\}$ heads and a position-wise feedforward network.  The
model dimension is $d_{\text{model}} \in \{64, 128\}$ and the
feedforward dimension is $d_{\text{ff}} = 4 \cdot d_{\text{model}}$.
Layer normalization is applied before each sub-layer (pre-norm
configuration).  A \emph{causal attention mask} is applied so that each
timestep attends only to current and past observations, ensuring the
model can operate in a streaming fashion at deployment.

\textbf{Output head.}  The transformer output at the final position
(timestep $t$) is passed through a small MLP with two output branches:
\begin{itemize}
\item Detection branch: linear layer $\to$ sigmoid, producing
  $p_t \in [0,1]$.
\item Magnitude branch: linear layer $\to$ softplus, producing
  $\hat{\delta}_t \geq 0$.
\end{itemize}

\textbf{Training.}  The transformer is trained in a supervised fashion on
data logged during the adversarial training phase
(Section~\ref{sec:data_logging}).  For each window
$(o_{t-L+1}, \ldots, o_t)$, the label consists of the ground-truth
disturbance indicator $y_t \in \{0,1\}$ (whether a disturbance was
applied at timestep $t$) and the ground-truth disturbance magnitude
$\delta_t = \|a_t^d\|_2$.  The training loss is
\begin{equation}\label{eq:detector_loss}
\cL = \cL_{\text{BCE}}(p_t, y_t) + \lambda\,\cL_{\text{MSE}}(\hat{\delta}_t, \delta_t),
\end{equation}
where $\cL_{\text{BCE}}$ is the binary cross-entropy loss,
$\cL_{\text{MSE}}$ is the mean squared error, and $\lambda > 0$ is a
weighting hyperparameter balancing the two objectives.

\textbf{Design for deployment.}  The transformer is intentionally kept
small to satisfy the real-time latency budget on Jetson hardware.  With
$N_L = 2$, $d_{\text{model}} = 64$, $N_H = 2$, and $L = 16$, the model
has approximately 50k parameters and can be executed in under 0.5\,ms
on a Jetson Orin after TensorRT optimization.

% ---------------------------------------------------------------------------
\subsection{Policy Mixing and Gating Rule}\label{sec:gating}
% ---------------------------------------------------------------------------

The gating function converts the transformer's outputs into the mixing
coefficient $\alpha_t$.  We consider two variants:

\textbf{Linear gating.}  A simple parameterization:
\begin{equation}\label{eq:gate_linear}
\alpha_t = \sigma(w_p \cdot p_t + w_\delta \cdot \hat{\delta}_t + b),
\end{equation}
where $\sigma(\cdot)$ is the logistic sigmoid, and $w_p, w_\delta, b$
are learnable scalar parameters.

\textbf{MLP gating.}  A more expressive variant:
\begin{equation}\label{eq:gate_mlp}
\alpha_t = f_\phi(p_t, \hat{\delta}_t),
\end{equation}
where $f_\phi$ is a small multi-layer perceptron (e.g., two hidden layers
of 32 units).  The MLP can learn nonlinear relationships such as
threshold effects or hysteresis-like behavior.

In both cases, the final action is computed as
\begin{equation}\label{eq:action_mix}
a_t = \alpha_t\,\pirob(s_t) + (1 - \alpha_t)\,\piopt(s_t).
\end{equation}

\textbf{Soft mixing vs.\ hard switching.}  An alternative design uses
hard switching: $\alpha_t \in \{0, 1\}$ based on a threshold on $p_t$.
While simpler, hard switching introduces \emph{chattering}---rapid
oscillation between policies near the decision boundary---which can
produce discontinuous control signals.  Soft mixing with a continuous
$\alpha_t$ avoids this problem and allows smooth transitions between the
two policies.  We include a comparison of soft and hard switching in the
evaluation plan (Section~\ref{sec:ablation}).

% ---------------------------------------------------------------------------
\subsection{Data Logging Specification}\label{sec:data_logging}
% ---------------------------------------------------------------------------

A standardized data-logging format is essential for reproducibility and
for generating the supervised training set for the transformer detector.
Table~\ref{tab:data_fields} specifies the fields logged at every timestep
during training.

\begin{table}[t]
\centering
\caption{Data logging specification.  Each row is recorded per timestep.}
\label{tab:data_fields}
\small
\begin{tabular}{@{}llp{4.0cm}@{}}
\toprule
\textbf{Field} & \textbf{Type} & \textbf{Description} \\
\midrule
\texttt{t} & \texttt{int} & Timestep index within episode \\
\texttt{episode\_id} & \texttt{int} & Unique episode identifier \\
\texttt{obs} & \texttt{float[$n$]} & Observation vector \\
\texttt{action\_ctrl} & \texttt{float[$m_c$]} & Controller action \\
\texttt{action\_dist} & \texttt{float[$m_d$]} & Disturbance/adversary action \\
\texttt{action\_total} & \texttt{float[$m_c$]} & Applied action after mixing \\
\texttt{reward} & \texttt{float} & Scalar reward \\
\texttt{terminated} & \texttt{bool} & Episode termination flag \\
\texttt{truncated} & \texttt{bool} & Episode truncation flag \\
\texttt{info/dist\_params} & \texttt{float[$k$]} & Ground-truth disturbance parameters \\
\texttt{disturbance\_tag} & \texttt{str} & Categorical disturbance label \\
\texttt{robust\_weight} & \texttt{float} & Target $\alpha_t$ for gating supervision \\
\bottomrule
\end{tabular}
\vspace{-0.3cm}
\end{table}

\textbf{Sequence window extraction.}  To train the transformer detector,
we extract sliding windows of length $L$ from the logged episodes.  For
timesteps $t < L$ within an episode, the window is left-padded with zeros
and an attention mask is constructed so that the transformer ignores
padded positions.  Windows \emph{never} cross episode boundaries: if the
remaining timesteps in an episode are fewer than $L$, the window is
zero-padded from the left rather than drawing observations from the next
episode.

\textbf{Storage format.}  Data is stored in HDF5 files (one per training
run) with chunked and compressed datasets, or in Apache Arrow/Parquet
format for efficient columnar access.  Each episode is a separate group
(HDF5) or partition (Arrow), enabling efficient random access during
window extraction.

% ---------------------------------------------------------------------------
\subsection{Algorithm Pseudocode}\label{sec:algorithms}
% ---------------------------------------------------------------------------

We present three algorithms covering the complete pipeline: adversarial
training, transformer detector training, and deployment inference.

\begin{algorithm}[t]
\caption{Off-Policy Adversarial Training Loop}
\label{alg:adversarial_training}
\begin{algorithmic}[1]
\renewcommand{\algorithmicrequire}{\textbf{Input:}}
\renewcommand{\algorithmicensure}{\textbf{Output:}}
\REQUIRE Environment $\text{env}$; actor networks $\piopt$, $\pirob$;
  critic $Q_\phi$; adversary $\piadv$; replay buffer $\mathcal{B}$;
  dataset logger $\cD$; episodes $N$; budget schedule $\epsilon(\cdot)$
\ENSURE Trained $\piopt$, $\pirob$, $\piadv$; dataset $\cD$
\STATE \textbf{Phase 1: Train optimal policy}
\FOR{episode $= 1$ \TO $N_{\text{opt}}$}
  \STATE $s_0 \leftarrow \text{env.reset}()$
  \FOR{$t = 0, 1, \ldots, T-1$}
    \STATE $a_t^c \leftarrow \piopt(s_t) + \text{noise}$
    \STATE $s_{t+1}, r_t, d_t \leftarrow \text{env.step}(a_t^c, \boldsymbol{0})$
    \STATE $\mathcal{B}.\text{add}(s_t, a_t^c, r_t, s_{t+1}, d_t)$
    \STATE $\cD.\text{log}(t, s_t, a_t^c, \boldsymbol{0}, r_t, d_t)$
    \STATE Update $\piopt$, $Q_\phi$ from $\mathcal{B}$ (SAC/TD3)
  \ENDFOR
\ENDFOR
\STATE \textbf{Phase 2: Train robust policy with adversary}
\FOR{episode $= 1$ \TO $N$}
  \STATE $\epsilon_t \leftarrow \epsilon(\text{episode})$ \hfill $\triangleright$ Curriculum
  \STATE $s_0 \leftarrow \text{env.reset}()$
  \FOR{$t = 0, 1, \ldots, T-1$}
    \STATE $a_t^c \leftarrow \pirob(s_t) + \text{noise}$
    \STATE $a_t^d \leftarrow \piadv(s_t)$, clipped to $\|a_t^d\| \leq \epsilon_t$
    \STATE $s_{t+1}, r_t, d_t \leftarrow \text{env.step}(a_t^c, a_t^d)$
    \STATE $\mathcal{B}.\text{add}(s_t, a_t^c, a_t^d, r_t, s_{t+1}, d_t)$
    \STATE $\cD.\text{log}(t, s_t, a_t^c, a_t^d, r_t, d_t)$
    \STATE Sample batch from $\mathcal{B}$
    \STATE Update $Q_\phi$: minimize TD error
    \STATE Update $\pirob$: maximize $Q_\phi(s, \pirob(s))$
    \STATE Update $\piadv$: minimize $Q_\phi(s, \pirob(s), \piadv(s))$
  \ENDFOR
\ENDFOR
\end{algorithmic}
\end{algorithm}

\begin{algorithm}[t]
\caption{Transformer Disturbance Detector Training}
\label{alg:transformer_training}
\begin{algorithmic}[1]
\renewcommand{\algorithmicrequire}{\textbf{Input:}}
\renewcommand{\algorithmicensure}{\textbf{Output:}}
\REQUIRE Dataset $\cD$ (logged episodes); window length $L$;
  transformer model $\cT_\theta$; learning rate $\eta$;
  loss weight $\lambda$; epochs $E$
\ENSURE Trained detector $\cT_\theta$
\STATE Extract sliding windows $\{(h_i, y_i, \delta_i)\}_{i=1}^{M}$
  from $\cD$
\STATE \quad $h_i = (o_{i-L+1}, \ldots, o_i)$, with zero-padding and
  attention masks
\STATE \quad $y_i = \mathbb{1}[\|a_i^d\|_2 > 0]$, \quad
  $\delta_i = \|a_i^d\|_2$
\STATE Split into training set $\cD_{\text{train}}$ and validation set
  $\cD_{\text{val}}$
\FOR{epoch $= 1$ \TO $E$}
  \FOR{each mini-batch $(H, \boldsymbol{y}, \boldsymbol{\delta})
    \subset \cD_{\text{train}}$}
    \STATE $(\hat{\boldsymbol{p}},\, \hat{\boldsymbol{\delta}})
      \leftarrow \cT_\theta(H)$ \hfill $\triangleright$ Forward pass
    \STATE $\cL \leftarrow \cL_{\text{BCE}}(\hat{\boldsymbol{p}},
      \boldsymbol{y}) + \lambda\,\cL_{\text{MSE}}(\hat{\boldsymbol{\delta}},
      \boldsymbol{\delta})$
    \STATE $\theta \leftarrow \theta - \eta\,\nabla_\theta \cL$
      \hfill $\triangleright$ Backpropagation
  \ENDFOR
  \STATE Evaluate on $\cD_{\text{val}}$; save checkpoint if improved
\ENDFOR
\RETURN $\cT_\theta$
\end{algorithmic}
\end{algorithm}

\begin{algorithm}[t]
\caption{Deployment Inference Loop (ROS~2 Node)}
\label{alg:deployment}
\begin{algorithmic}[1]
\renewcommand{\algorithmicrequire}{\textbf{Input:}}
\REQUIRE TensorRT engines: $\piopt$, $\pirob$, $\cT_\theta$;
  gating parameters $\phi$; history buffer $H$; control frequency $f$
\STATE Initialize ROS~2 node, sensor subscriber, action publisher
\STATE Initialize CUDA streams for pipelined inference
\STATE Set timer period $\Delta t \leftarrow 1/f$
\STATE \textbf{On timer callback:}
\STATE \quad $o_t \leftarrow$ read latest sensor message
\STATE \quad Append $o_t$ to history buffer $H$
\IF{$|H| < L$}
  \STATE \quad Construct zero-padded input with attention mask
\ENDIF
\STATE \quad $(\hat{p}_t,\, \hat{\delta}_t) \leftarrow
  \cT_\theta(H[-L:])$ \hfill $\triangleright$ TensorRT
\STATE \quad $\alpha_t \leftarrow f_\phi(\hat{p}_t, \hat{\delta}_t)$
  \hfill $\triangleright$ Gating
\STATE \quad $a^{\text{opt}}_t \leftarrow \piopt(o_t)$
  \hfill $\triangleright$ TensorRT
\STATE \quad $a^{\text{rob}}_t \leftarrow \pirob(o_t)$
  \hfill $\triangleright$ TensorRT
\STATE \quad $a_t \leftarrow \alpha_t\,a^{\text{rob}}_t +
  (1 - \alpha_t)\,a^{\text{opt}}_t$
\STATE \quad Publish $a_t$ to actuator topic
\STATE \quad Log: timestamp, $\alpha_t$, $\hat{p}_t$, $\hat{\delta}_t$,
  latency
\end{algorithmic}
\end{algorithm}

Algorithm~\ref{alg:adversarial_training} describes the two-phase off-policy
training loop.  Phase~1 trains the optimal policy without any adversary.
Phase~2 trains the robust policy against the adversary, with curriculum
scheduling of the disturbance budget.  Both phases log all transitions
to the dataset $\cD$.

Algorithm~\ref{alg:transformer_training} describes the supervised training
of the transformer disturbance detector using the logged data.

Algorithm~\ref{alg:deployment} describes the real-time inference loop
running as a ROS~2 node on the Jetson platform.

% ===========================================================================
\section{System Implementation}\label{sec:implementation}
% ===========================================================================

% ---------------------------------------------------------------------------
\subsection{Buffer and Logger Utilities}\label{sec:buffers}
% ---------------------------------------------------------------------------

The system uses two distinct buffer types depending on the training
algorithm.

\textbf{Replay buffer (off-policy).}  A fixed-capacity circular buffer
stores transitions $(s, a^c, a^d, r, s', d)$.  The buffer supports
uniform random sampling and, optionally, Prioritized Experience Replay
(PER), which samples transitions with probability proportional to their
TD error magnitude.  The buffer capacity is typically set to $10^6$
transitions.

\textbf{Rollout buffer (on-policy).}  For PPO, a rollout buffer stores
complete episodes (or fixed-length trajectory segments) including
observations, actions, rewards, log-probabilities, and value estimates.
At the end of each rollout, Generalized Advantage Estimation (GAE) is
computed with parameters $\lambda_{\text{GAE}} = 0.95$ and $\gamma = 0.99$.
The buffer is emptied after each policy update.

\textbf{Dataset logger.}  Independent of the replay or rollout buffer, the
dataset logger writes every transition to persistent storage in the format
specified in Table~\ref{tab:data_fields}.  This logger is active during
both Phase~1 (optimal policy training) and Phase~2 (adversarial training),
producing the complete dataset required for transformer training.  The
logger writes in append-only mode with periodic flushing to minimize I/O
overhead during training.

% ---------------------------------------------------------------------------
\subsection{Training Pipeline}\label{sec:training_pipeline}
% ---------------------------------------------------------------------------

The training pipeline consists of four phases:

\textbf{Phase~1: Optimal policy ($\piopt$).}  Train using SAC or TD3 on
the nominal Gymnasium/MuJoCo environment.  Hyperparameters: learning rate
$3 \times 10^{-4}$, batch size 256, replay buffer size $10^6$, target
smoothing coefficient $\tau = 0.005$.  Training typically requires
$1$--$3 \times 10^6$ environment steps depending on the task complexity.

\textbf{Phase~2: Robust policy ($\pirob$) and adversary ($\piadv$).}
Train in the zero-sum game using the modified environment that accepts
both controller and adversary actions.  The adversary's disturbance
budget is annealed linearly from $0$ to $\epsilon_{\max}$ over the first
$20\%$ of training.  Both the controller and adversary use separate
optimizers but share the same replay buffer.  The critic is updated with
every environment step; the controller and adversary actors are updated
every two environment steps (delayed update as in TD3).

\textbf{Phase~3: Transformer detector ($\cT_\theta$).}  Extract sliding
windows from the logged dataset $\cD$.  Train the transformer with
AdamW optimizer, learning rate $10^{-4}$, weight decay $10^{-2}$, batch
size 128, for 100 epochs.  Use cosine learning rate scheduling with
warmup over the first 5 epochs.  The loss weight is set to $\lambda = 1.0$
and tuned on a held-out validation set.

\textbf{Phase~4: Gating network (optional).}  If the linear gating
rule~\eqref{eq:gate_linear} is insufficient, train the MLP gating
network~\eqref{eq:gate_mlp} on a separate validation set of episodes
where the target $\alpha_t$ is computed from the ground-truth disturbance
labels.  Alternatively, the gating parameters can be tuned by optimizing
the mixed policy's return in simulation.

\textbf{Migration to JAX/MJX.}  The current PyTorch implementation
processes one environment instance per training step.  A planned migration
to JAX with MJX~\cite{zakka2025mujoco} will enable
GPU-vectorized simulation of thousands of environment instances in
parallel, as demonstrated by Brax~\cite{freeman2021brax}.  This is
expected to yield speedups on the order of $1000\times$ for on-policy
methods and $10$--$100\times$ for off-policy methods (where the
bottleneck shifts from data collection to gradient computation).  The
JAX implementation will use Flax for neural network parameterization and
Optax for optimization.

% ---------------------------------------------------------------------------
\subsection{Export Pipeline}\label{sec:export}
% ---------------------------------------------------------------------------

All neural network components must be converted from their training
framework representation to TensorRT engines for deployment.

\textbf{PyTorch export path.}  Models are first traced using
\texttt{torch.jit.trace} with representative input tensors, then exported
to ONNX format via \texttt{torch.onnx.export} with opset version 17.
For the policy networks, the input is a single observation vector
$o \in \R^d$ and the output is an action vector $a \in \R^{m_c}$
(static shapes).  For the transformer, the input is a window
$H \in \R^{L \times d}$ with a fixed sequence length $L$ (static shape).

\textbf{JAX export path.}  JAX models are exported using
\texttt{jax.export} or the \texttt{jax2tf} bridge to produce a
TensorFlow SavedModel, which is then converted to ONNX using
\texttt{tf2onnx}.  Alternatively, the \texttt{jax2onnx} converter can be
used for direct export.

\textbf{ONNX to TensorRT.}  The ONNX models are compiled to TensorRT
engines using \texttt{trtexec} or the TensorRT Python
API~\cite{nvidia2024tensorrt}.  Key optimization settings include:
\begin{itemize}
\item \textbf{Precision:} FP16 mode is used by default, providing a
  $2\times$ speedup over FP32 with negligible accuracy loss for the
  policy networks.  INT8 quantization is explored for the transformer,
  using calibration data from the training
  set~\cite{onnxruntime2022transformer}.
\item \textbf{Layer fusion:} TensorRT automatically fuses
  convolution+bias+activation, attention QKV projections, and other
  patterns.
\item \textbf{Engine caching:} TensorRT engines are serialized to disk
  and reloaded on subsequent runs, avoiding the overhead of re-optimization
  at startup~\cite{jeon2022tensorrt}.
\end{itemize}

For the transformer, the causal attention mask is a static boolean tensor
that can be baked into the TensorRT engine as a constant, eliminating
the need to pass it as a dynamic input.

% ---------------------------------------------------------------------------
\subsection{ROS~2 and Isaac ROS Integration}\label{sec:ros2}
% ---------------------------------------------------------------------------

The deployment system is implemented as a ROS~2 (Humble or later) node
with the following structure:

\textbf{Subscriptions.}  The node subscribes to sensor topics providing
joint states, IMU data, and/or other proprioceptive observations.
Messages are received via a quality-of-service (QoS) profile configured
for best-effort reliability and volatile durability, minimizing latency.

\textbf{Timer callback.}  A timer triggers the inference loop
(Algorithm~\ref{alg:deployment}) at the configured control frequency.
The callback is assigned to a high-priority callback group within a
\texttt{SingleThreadedExecutor}.

\textbf{Publication.}  The computed action $a_t$ is published to an
actuator command topic.  Messages are pre-allocated at initialization to
avoid dynamic memory allocation in the callback.

\textbf{NITROS integration.}  When available, Isaac ROS
NITROS~\cite{nvidia2024isaacros} is used for zero-copy GPU tensor
transfer between the sensor preprocessing node and the inference node,
eliminating CPU-GPU memory copies.

\textbf{Real-time considerations.}  The callback avoids all dynamic
memory allocation, uses pre-allocated CUDA memory for TensorRT input/output
buffers, and employs intra-process communication when the sensor node and
inference node run in the same process.  The CUDA context is initialized
once at startup, and TensorRT execution contexts are created and
retained for the lifetime of the node.

% ---------------------------------------------------------------------------
\subsection{Latency Budgeting and Scheduling}\label{sec:latency}
% ---------------------------------------------------------------------------

Table~\ref{tab:latency} presents the target latency budget for each
component of the inference pipeline at a control frequency of 500\,Hz
(2\,ms total budget) on an NVIDIA Jetson Orin in MAXN power mode.

\begin{table}[t]
\centering
\caption{Inference latency budget at 500\,Hz on Jetson Orin (MAXN).}
\label{tab:latency}
\small
\begin{tabular}{@{}lr@{}}
\toprule
\textbf{Component} & \textbf{Target latency} \\
\midrule
Transformer forward pass ($\cT_\theta$)
  & $< 0.50$\,ms \\
Policy forward pass ($\piopt$)
  & $< 0.30$\,ms \\
Policy forward pass ($\pirob$)
  & $< 0.30$\,ms \\
Gating computation ($f_\phi$)
  & $< 0.05$\,ms \\
ROS message serialization + publish
  & $< 0.20$\,ms \\
\midrule
\textbf{Total}
  & $< 1.35$\,ms \\
\bottomrule
\end{tabular}
\vspace{-0.1cm}
\end{table}

The total target of 1.35\,ms leaves a margin of 0.65\,ms relative to
the 2\,ms budget, accommodating occasional latency spikes.  Several
strategies are employed to further reduce latency:

\textbf{CUDA streams.}  The two policy forward passes are independent and
can be executed concurrently on separate CUDA streams.  On the Jetson
Orin's GPU (which supports concurrent kernel execution), this reduces the
effective policy inference time from $2 \times 0.30 = 0.60$\,ms to
approximately $0.35$\,ms.

\textbf{Pipelining.}  The transformer forward pass and policy forward
passes are partially overlapped: the transformer can begin processing
while the previous timestep's action is being published.  This is
implemented by maintaining a one-step pipeline depth in the CUDA
execution schedule.

\textbf{Power modes.}  The Jetson Orin supports multiple power modes
ranging from 15\,W to MAXN ($\sim$60\,W).  In MAXN mode, all GPU and
CPU cores are active at maximum frequency, providing the lowest latency.
For battery-powered applications, the 30\,W mode is a practical compromise,
with an expected latency increase of approximately $1.5$--$2\times$.

% ===========================================================================
\section{Evaluation Plan}\label{sec:evaluation}
% ===========================================================================

This section describes the planned experimental evaluation.  No empirical
results are reported; the purpose is to specify the evaluation protocol
in sufficient detail for reproducibility.

% ---------------------------------------------------------------------------
\subsection{Benchmark Environments}\label{sec:benchmarks}
% ---------------------------------------------------------------------------

We plan to evaluate on four MuJoCo locomotion environments available
through Gymnasium~\cite{towers2024gymnasium}:
\begin{itemize}
\item \textbf{HalfCheetah-v4}: 17-dimensional observation, 6-dimensional
  action.  A planar biped that must run forward.
\item \textbf{Walker2d-v4}: 17-dimensional observation, 6-dimensional
  action.  A planar biped that must walk while maintaining balance.
\item \textbf{Ant-v4}: 27-dimensional observation, 8-dimensional action.
  A 3D quadruped navigating in the $xy$-plane.
\item \textbf{Humanoid-v4}: 376-dimensional observation, 17-dimensional
  action.  A 3D bipedal humanoid, the most challenging benchmark.
\end{itemize}

For each environment, disturbances are introduced through three mechanisms:
(i)~random external forces applied to the torso at random intervals,
(ii)~perturbation of physical parameters (mass, friction, damping) by up
to $\pm30\%$ from nominal values, and (iii)~additive Gaussian noise on
observations with standard deviation up to $0.1$.

% ---------------------------------------------------------------------------
\subsection{Metrics}\label{sec:metrics}
% ---------------------------------------------------------------------------

The evaluation uses the following metrics:
\begin{itemize}
\item \textbf{Nominal return} $J_{\text{nom}}$: expected discounted return
  in the nominal environment (no disturbance).
\item \textbf{Robust return} $J_{\text{rob}}$: expected discounted return
  under the worst-case learned adversary at budget $\epsilon_{\max}$.
\item \textbf{Worst-case regret}:
  $\Delta J = J_{\text{nom}}^* - J_{\text{rob}}$, where
  $J_{\text{nom}}^*$ is the return of the best nominal-only policy.
  Lower is better.
\item \textbf{Constraint violation rate}: fraction of timesteps in which
  joint limits, self-collisions, or other safety constraints are violated.
\item \textbf{Recovery time}: number of timesteps from disturbance onset
  to return within $95\%$ of nominal tracking performance.
\item \textbf{Control frequency}: achieved Hz on Jetson hardware.
\item \textbf{Inference latency}: mean, 95th percentile (p95), and 99th
  percentile (p99) of the end-to-end inference time.
\item \textbf{Jitter}: standard deviation of inter-callback timing,
  measuring real-time determinism.
\end{itemize}

% ---------------------------------------------------------------------------
\subsection{Ablation Studies}\label{sec:ablation}
% ---------------------------------------------------------------------------

The following ablation studies are planned to isolate the contribution of
each component:
\begin{enumerate}
\item \textbf{No adversary}: use $\piopt$ only (i.e., standard RL without
  robustness).  This establishes the nominal performance ceiling and
  measures degradation under disturbance.
\item \textbf{No transformer detector}: fix $\alpha_t$ to a constant value
  (e.g., $\alpha = 0.5$) rather than using the learned gating.  This
  measures the value of runtime disturbance awareness.
\item \textbf{Hard switching vs.\ soft mixing}: replace the continuous
  $\alpha_t$ with a binary decision
  $\alpha_t \in \{0, 1\}$ based on a threshold on $p_t$.  This measures
  the benefit of smooth blending.
\item \textbf{History window length}: vary $L \in \{8, 16, 32, 64\}$ to
  assess the trade-off between detection accuracy and computational cost.
\item \textbf{Adversary budget}: vary $\epsilon$ over a range to
  characterize the robustness-performance Pareto frontier.
\item \textbf{Single adversary vs.\ adversarial population}: compare
  training against a single adversary with training against a population
  of diverse adversaries~\cite{vinitsky2020robust}, measuring
  generalization to unseen disturbance patterns.
\end{enumerate}

% ---------------------------------------------------------------------------
\subsection{Baselines}\label{sec:baselines}
% ---------------------------------------------------------------------------

We compare against the following baselines:
\begin{itemize}
\item \textbf{SAC/TD3 (nominal)}: standard deep RL without any robustness
  mechanism, representing the performance ceiling under nominal conditions
  and the floor under disturbances.
\item \textbf{RARL}~\cite{pinto2017robust}: single robust policy trained
  with a co-learned adversary applying body forces.  No gating or runtime
  detection.
\item \textbf{SA-MDP}~\cite{zhang2020robust}: single robust policy trained
  with state-adversarial perturbations bounded in $\ell_\infty$ norm.
  No gating.
\item \textbf{Domain randomization (DR)}~\cite{tobin2017domain}: single
  policy trained with randomized physical parameters (mass, friction,
  damping) drawn uniformly from a specified range.  No adversary, no
  gating.
\item \textbf{Ours (full)}: dual-policy architecture with transformer
  detector and learned gating, as described in this paper.
\end{itemize}

All methods use the same neural network architecture (two hidden layers,
256 units, ReLU) for the policy networks, and are trained for the same
total number of environment interactions to ensure a fair comparison.

% ---------------------------------------------------------------------------
\subsection{Safety Checks and Failure Cases}\label{sec:safety}
% ---------------------------------------------------------------------------

We define a \emph{critical failure} as any of the following events:
(i)~joint position exceeds its mechanical limit, (ii)~self-collision
between non-adjacent links, (iii)~the robot falls (torso height below a
task-specific threshold).  These events are monitored across all
evaluation runs, and the critical failure rate is reported alongside the
return metrics.

Planned stress tests include:
\begin{itemize}
\item \textbf{Out-of-distribution adversary}: apply disturbances with
  magnitude $\delta > \epsilon_{\max}$ (beyond the training budget) and
  measure graceful degradation.
\item \textbf{Sensor dropout}: randomly zero out observation dimensions
  with probability $p_{\text{drop}} \in \{0.05, 0.1, 0.2\}$ to simulate
  sensor failures.
\item \textbf{Communication delay}: inject artificial latency of 1--5\,ms
  into the observation pipeline to simulate network delays in distributed
  systems.
\item \textbf{Adversary-fools-detector}: construct adversarial
  disturbances that are specifically designed to evade the transformer
  detector (low $p_t$ despite active disturbance), and measure the
  resulting performance degradation.
\end{itemize}

% ===========================================================================
\section{Discussion and Limitations}\label{sec:discussion}
% ===========================================================================

While the proposed methodology provides a principled framework for robust
real-time control, several limitations and open questions merit discussion.

\textbf{Detector evasion.}  The transformer disturbance detector is
trained on data from the adversary encountered during training.  A
sufficiently sophisticated adversary could learn to apply disturbances
that evade detection---perturbations that degrade performance without
producing observable signatures in the observation history.  This
represents a fundamental limitation of any detection-based gating scheme
and motivates future work on adversarially robust detectors.

\textbf{Sim-to-real gap for the detector.}  The transformer is trained
entirely on simulated data.  The observation statistics on a physical
robot may differ from simulation due to sensor noise, calibration errors,
and unmodeled dynamics.  Domain randomization of sensor characteristics
during data generation may partially address this issue, but a gap is
likely to remain.  Fine-tuning the detector on a small amount of real
data is a practical mitigation strategy.

\textbf{Compute constraints on Jetson.}  The Jetson Orin's GPU, while
capable, imposes tight constraints on transformer size.  The current
design targets 2--4 layers with 64--128 model dimension, which limits
the complexity of temporal patterns the detector can capture.
Larger models may be feasible on future embedded GPU platforms or through
more aggressive quantization.

\textbf{Training cost.}  The three-phase training pipeline (optimal
policy, robust policy with adversary, transformer detector) requires
approximately $3\times$ the compute of standard RL training.  The
migration to JAX/MJX with GPU-vectorized environments is expected to
substantially reduce wall-clock time, but the fundamental sample
complexity remains higher due to the adversarial game dynamics.

\textbf{Known disturbance bounds.}  The method assumes a known upper
bound $\epsilon$ on the disturbance magnitude.  In practice, this bound
must be estimated from domain knowledge or system identification.  If
the true disturbance exceeds $\epsilon_{\max}$, the robust policy's
guarantees no longer hold, though the dual-policy architecture may still
provide better performance than a nominal-only policy.

\textbf{Action-space mixing.}  The policy mixing rule~\eqref{eq:action_mix}
operates in the action space, which may not be optimal.  If the two
policies produce actions in different regions of the action space, a
convex combination may yield an intermediate action that neither policy
would choose.  An alternative is to mix in a \emph{latent space} shared
by the two policies, or to use a hypernetwork that generates policy
parameters conditioned on $\alpha_t$.  We leave exploration of these
alternatives to future work.

\textbf{Scalability to high-dimensional observations.}  The current
formulation assumes proprioceptive observations of moderate dimension
($d \leq 400$).  Extension to high-dimensional observations (e.g.,
images from onboard cameras) would require replacing the transformer
input with learned feature vectors from a vision encoder, significantly
increasing the computational requirements and the complexity of the
deployment pipeline.

% ===========================================================================
\section{Conclusion and Future Work}\label{sec:conclusion}
% ===========================================================================

We have presented a methodology for robust real-time control based on
zero-sum Markov games with a dual-policy mixing architecture.  The
approach maintains two separately trained policies---one optimal for
nominal conditions, one robust to worst-case adversarial
disturbances---and uses a lightweight transformer encoder to detect
disturbances from observation history and continuously blend the two
policies via a learned gating function.  The complete system includes
a training pipeline with standardized data logging, an export pipeline
through ONNX to TensorRT for embedded GPU deployment, and a ROS~2
integration designed to meet sub-two-millisecond latency budgets on the
NVIDIA Jetson platform.  A comprehensive evaluation plan with ablation
studies, baselines, and safety checks has been specified.

Future work will proceed along several directions:
\begin{enumerate}
\item \textbf{End-to-end training.}  Replace the three-phase sequential
  pipeline with a single end-to-end training procedure that jointly
  optimizes the two policies, the transformer detector, and the gating
  network through a unified objective.
\item \textbf{Vision-based disturbance detection.}  Extend the transformer
  detector to process visual observations (e.g., depth images) in addition
  to proprioceptive data, enabling detection of visually observable
  disturbances such as approaching obstacles or terrain changes.
\item \textbf{Multi-agent extension.}  Generalize the framework to
  cooperative and competitive multi-agent settings, where multiple
  controllers must coordinate their robust strategies against a shared
  or distributed adversary.
\item \textbf{Formal safety guarantees.}  Integrate Control Barrier
  Functions (CBFs) into the gating mechanism to provide formal safety
  certificates, ensuring that the mixed action always satisfies safety
  constraints regardless of the detector's output.
\item \textbf{Empirical validation.}  Execute the evaluation plan
  described in Section~\ref{sec:evaluation} across all benchmark
  environments and on physical robot hardware, reporting the quantitative
  results in a follow-up study.
\end{enumerate}

% ===========================================================================
\section*{Acknowledgements}
% ===========================================================================

This work was supported by [Funding Agency, Grant Number].  The authors
thank [Collaborators] for helpful discussions.

% ===========================================================================
% Bibliography
% ===========================================================================
\bibliographystyle{IEEEtran}
\bibliography{refs}

\end{document}
